\documentclass[11pt]{article}
\usepackage[margin=1in]{geometry}
\usepackage{amsmath,amssymb}
\usepackage{booktabs}
\usepackage{longtable}
\usepackage{hyperref}
\hypersetup{colorlinks=true,linkcolor=blue,urlcolor=blue}
\usepackage{fancyvrb}

\newcommand{\lean}[1]{\texttt{#1}}

\title{A Lean 4 formalization of\\ ``$\zeta(5)$ is irrational'' (A.\ Fauzan, 17 September 2026)}
\author{Summary of \url{https://github.com/danromik/zeta5-irrationality}}
\date{24 September 2026}

\begin{document}
\maketitle
\sloppy

\begin{abstract}
The repository contains a complete formalization in Lean~4, with Mathlib, of the proof in
A.~Fauzan's preprint ``$\zeta(5)$ is irrational''. The main theorem,
\lean{Zeta5.zeta5\_irrational : Irrational zeta5}, is proved with no \lean{sorry}. It depends on
the three standard axioms of Lean and on one external axiom, the prime number theorem in the
form $\theta(x)\sim x$, which Mathlib does not contain. The formal proof passes through the
paper's intermediate statements (Theorem~2.1, Propositions~2.2, 4.1, 4.3, 5.1, 5.2, 6.3,
Lemma~3.3 and the numbered equations between them); at the places listed in
Section~\ref{sec:dev} it proves a step by a different argument from the paper's.
\end{abstract}

\section{The theorem and its assumptions}

\begin{Verbatim}[fontsize=\small,frame=single,commandchars=\\\{\}]
def zeta5 : \(\mathbb{R}\) := \(\sum\)' v : \(\mathbb{N}\), (1 : \(\mathbb{R}\)) / ((v : \(\mathbb{R}\)) + 1) ^ 5
theorem zeta5_irrational : Irrational zeta5
\end{Verbatim}
\noindent Both are in the namespace \lean{Zeta5}; \lean{Irrational} is Mathlib's predicate.
 \lean{\#print axioms Zeta5.zeta5\_irrational} gives
\begin{Verbatim}[fontsize=\small,frame=single]
[propext, Classical.choice, Quot.sound, Zeta5.Axioms.chebyshev_theta_asymptotic]
\end{Verbatim}
\noindent The last is the only axiom declared in the project (\lean{Zeta5/Axioms.lean}):
\begin{Verbatim}[fontsize=\small,frame=single,commandchars=\\\{\}]
axiom chebyshev_theta_asymptotic :
    Tendsto (fun x : \(\mathbb{R}\) => Chebyshev.theta x / x) atTop (nhds 1)
\end{Verbatim}
\noindent where \lean{Chebyshev.theta} is Mathlib's $\theta(x)=\sum_{p\le x}\log p$. This is the
prime number theorem (Hadamard, de~la~Vall\'ee~Poussin, 1896) as stated in textbooks, e.g.\
Apostol, \emph{Introduction to Analytic Number Theory}, Thm.~4.4 with Ch.~13; the paper cites
DLMF 27.12.2--27.12.4. It is used only in Proposition~5.2, through the partial-summation
theorem \lean{Zeta5.PNT.prime\_riemann\_sum}. Lean therefore checks the implication
\[
  \text{prime number theorem} \;\Longrightarrow\; \zeta(5)\notin\mathbb{Q}.
\]

\section{The build}

Lean 4.34.0, Mathlib v4.34.0; 70 files under \lean{Zeta5/}, about 29\,400 lines.
\lean{lake build} ends with \lean{Build completed successfully (8995 jobs)}, with no error and
no \lean{sorry} warning; the 17 warnings are linter warnings (unused section variables and
unreferenced hypothesis names). A clean build takes about 7 minutes on a 12-core machine.

\section{Structure}

\begin{longtable}{p{0.27\textwidth}p{0.66\textwidth}}
\toprule
Paper & Lean files (in \lean{Zeta5/}) \\
\midrule
\endhead
\S2, Proposition 2.2 & \lean{Functional}, \lean{Positivity}; the pole integral of p.~5 in
\lean{Hermite} \\
\S3, Lemma 3.3, (3.11)--(3.12) & \lean{LocalFunctional}, \lean{Lemma33}, \lean{CrudeBound},
\lean{Arithmetic} \\
\S4.1, Proposition 4.1 & \lean{Counting}, \lean{Section41}, \lean{Section3}, \lean{InnerTate},
\lean{InnerGeneral}, \lean{InnerEntries}, \lean{HermiteBasis} \\
\S4.2, Proposition 4.3 & \lean{Lemma42}, \lean{OuterLocal}, \lean{OuterBasis},
\lean{OuterRange} \\
\S5, Propositions 5.1, 5.2 & \lean{Normalization}, \lean{Uniformity}, \lean{PNT},
\lean{PrimeSum}, \lean{Tail} \\
Appendix B & \lean{AppendixB} \\
\S6, Appendix A, Proposition 6.3 & \lean{RealBound}, \lean{Sec6/*} (37 files) \\
\S7, Theorems 2.1, 1.1 & \lean{Asymptotics}, \lean{Skeleton}, \lean{Interface} \\
\bottomrule
\end{longtable}

\noindent \lean{Interface.lean} contains, for each numbered statement of the paper used by
Theorem~1.1, a declaration whose type transcribes the printed statement. The finite
computations of the paper (Appendix~B's 143 intervals and Tables~3--4, Table~1 and the
constants of Appendix~A.4, the margin (7.2)) and the certified partition used for (6.2) are
proved by exact rational arithmetic and certified enclosures; the project contains no
\lean{native\_decide}, \lean{implemented\_by} or \lean{extern}.

\section{Deviations from the paper}\label{sec:dev}

Each item is recorded in the docstring of the declaration concerned.
\begin{enumerate}
\item \textbf{Lemma 3.2, (3.7).} The Tate algebra, (3.5) and (3.6) are not formalized; the
valuation bound that Proposition~4.1 needs is proved from Raabe's multiplication theorem for
$\tau$ and a truncated inverse of the far-pole factor.
\item \textbf{Lemma 3.3, polynomial part.} Proved through the binomial basis $\binom{x+K}{k}$
and an exact computation of $\tau$, in place of (3.8)--(3.9). The term $-v_p(24)$ of (3.10) is
unnecessary.
\item \textbf{Unimodularity of (4.5) and (4.11)}, from one lemma on bases that reduce mod $p$ to
a Hermite-interpolation basis.
\item \textbf{The inequality $b_a\le 6\alpha x+3$ (p.~10)}, stated in the paper without proof,
is proved in the strict form $p\,\ell_N(a) < 2N+p$.
\item \textbf{The divided difference of p.~12} uses the term $-1/4+1/(2j)$ of (2.3) in addition
to the cited congruence.
\item \textbf{(4.10)}: the vanishing is proved in the monomial basis, and
$\operatorname{rank}L\le r_p$ is transported to the basis (4.11).
\item \textbf{(5.16)--(5.17)} are proved in the combined form that (5.18) uses.
\item \textbf{The pole integral of p.~5} is proved from the series for $w$ and the
Mittag-Leffler expansion of the cotangent, not from Hermite's formula.
\item \textbf{Partial summation (p.~15)} is proved from $\theta(x)\sim x$ by a Darboux-type
sandwich.
\item \textbf{(6.6)--(6.9).} The kernel is regularized on the real line,
$\tfrac12\log(x^2+\varepsilon^2)$ with $\varepsilon=(16K)^{-2}$, in place of circles in
$\mathbb{C}$; Lemma~6.2 is proved for this kernel from a Frullani representation and the
positive definiteness of the Gaussian kernel. The smoothing error is bounded directly. The
configuration bound obtained has $20h$ in place of $(120+\sqrt2)h$.
\item \textbf{(A.2)} is used as the inequality $I(\rho)\le I_k(\rho,\rho)$ for the regularized
kernel.
\item \textbf{(6.10)--(6.13).} The factor $1/h!$ is kept, and $326K^6$ replaces $652$.
\item \textbf{(6.2) on $[0,2]$} is proved by kernel-checked interval arithmetic on the closed
forms (A.1), (A.5) over a partition of 1049 cells, in place of (A.9) on Table~2; for $t\ge2$,
two cells cover $[2,4]$ and an analytic bound covers $t\ge4$.
\end{enumerate}

\section{Verification}

To accept the result, a reader needs to trust Lean's kernel and to read two statements, both
written with Mathlib names only in \lean{cert/Check.lean}:
\begin{Verbatim}[fontsize=\small,frame=single,commandchars=\\\{\}]
example : Irrational (\(\sum\)' v : \(\mathbb{N}\), (1 : \(\mathbb{R}\)) / ((v : \(\mathbb{R}\)) + 1) ^ 5) :=
  Zeta5.zeta5_irrational
example : Tendsto (fun x : \(\mathbb{R}\) => Chebyshev.theta x / x) atTop (nhds 1) :=
  Zeta5.Axioms.chebyshev_theta_asymptotic
\end{Verbatim}
\noindent together with the Mathlib definitions of \lean{Irrational}, \lean{tsum} and
\lean{Chebyshev.theta}. The checks are \lean{lake build}, \lean{lake env lean cert/Check.lean}
(which prints the axiom list of Section~1), and \lean{lake env leanchecker Zeta5}, which
re-checks every declaration of the compiled modules with the kernel alone
(\lean{docs/CERTIFICATION.md}).

\end{document}
